\section{A sharp height bound for fixed-port phases}
\label{sec:height}

\subsection{The algebra of active subpartitions}

The effect of a cone sequence on the source atoms has a compact canonical
description. Record the active atom set $A$ and partition it into blocks,
where two active atoms lie in the same block when the applied cone groups
connect them. Each block means that all points in the union of its atoms
are identified. Atoms outside $A$ impose no new relation. This is an
\emph{active subpartition}, not just an ordinary partition of all $m$ atoms.
In particular, an active singleton block must remain visible.

The join $S\vee T$ of two active subpartitions has active set the union of
their active sets and equivalence blocks generated by both collections of
blocks. It is computed by ordinary finite set union. The empty collection
is the identity.

\begin{lemma}[Commutative idempotent action]
\label{lem:semilattice}
Joining summaries agrees with composing their cone actions. Thus
\[
 (S\vee T)\vee U=S\vee(T\vee U),\qquad
 S\vee T=T\vee S,\qquad S\vee S=S.
\]
These identities hold for every initial relation $E$, not only for the
initial discrete relation.
\end{lemma}

\begin{proof}
Let $R_S$ be the union of the full relations on the point sets represented
by the blocks of $S$. Applying $S$ and then $T$ gives the equivalence
closure of $E\cup R_S\cup R_T$. Its closure depends neither on the order
nor on repeated copies of these generating relations. Two atom groups
intersect precisely when their point unions share an atom, so the connected
blocks of the join generate the same full relations. Associativity,
commutativity and idempotence follow.
\end{proof}

For example, coning atom $0$, then atom $1$, does \emph{not} mean that atoms
$0$ and $1$ are connected: the summary has two singleton blocks. Coning
$\{0,1\}$ instead has one two-element block. Dropping singleton blocks
would incorrectly erase the first two operations.

\begin{theorem}[Sharp fixed-atom height]
\label{thm:height}
For $m\ge1$, a sequence of full-port cones on $m$ fixed nonempty atoms
contains at most $2m-1$ steps at which the component partition changes.
Equivalently, its component count has at most $2m-1$ strict drops.
The bound is attained for every $m$. For $m=0$ there are no changes.
\end{theorem}

\begin{proof}
For an active subpartition with $a$ active atoms and $b$ blocks, set
\[
 \Phi=2a-b.
\]
The empty summary has $\Phi=0$. A nonempty summary has
$0\le\Phi\le2m-1$, since $a\le m$ and $b\ge1$.

Suppose a new nonempty cone group contains $k$ previously inactive atoms
and meets $q$ old active blocks. Its update replaces those $q$ blocks by
one block also containing the $k$ new atoms. Consequently
\[
 \Delta\Phi=2k+q-1.
\]
When the summary really changes, either $k\ge1$, or $k=0$ and $q\ge2$;
in both cases $\Delta\Phi\ge1$. If $k=0$ and $q=1$, the summary is
unchanged. An empty cone is also unchanged.

By \cref{lem:semilattice}, an unchanged summary cannot change the relation
on points, whatever the initial $E$. Every actual partition change therefore
uses at least one unit of the bounded potential, giving the result.
A strict coarsening of a finite partition decreases its number of blocks,
so partition changes and strict count drops are equivalent here.

For sharpness, take $2m$ initially isolated points and $m$ atoms of size two.
Cone each singleton atom in turn, producing $m$ strict drops. Then cone
$\{0,j\}$ for $j=1,\ldots,m-1$, producing another $m-1$ strict drops.
There are exactly $2m-1$ changes.
\end{proof}

This potential is qualitatively different from a large decreasing binary
counter. It has height linear in the number of fixed atoms, independent of
their capacities, component multiplicities, and the number of redundant
instructions. The hypothesis ``fixed'' cannot be omitted. An operation that
introduces a new atom or removes an old identification changes the state
space in which the potential is measured.

\section{Compressed programs, threshold search and complete event histories}
\label{sec:programs}

\subsection{Straight-line program contract}

A program is an acyclic list of $g$ nodes. A node is one of
\[
 \Cone(A),\qquad \operatorname{concat}(P,Q),\qquad
 \operatorname{power}(P,p),\quad p\in\N.
\]
References point strictly to earlier nodes. A root selects the requested
program. A cone is one instruction, including an empty cone. Concatenation
adds lengths and power multiplies a length by $p$. A zero power is an empty
program. Exponents and lengths are exact binary integers.

This language contains fixed atom names in its leaves. It does not contain
an instruction of the form ``shift this port by the current iteration
number'', an unknown attachment selected adaptively from new geometry, or
a relabeling hidden inside a power. Such languages require additional
semantics and are outside the following proofs.

Let $L(P)$ be expanded length and $S(P)$ the active-subpartition summary.
Bottom-up compilation uses
\begin{align*}
 L(\Cone(A))&=1,&S(\Cone(A))&=\{A\}\quad(A\ne\varnothing),\\
 L(PQ)&=L(P)+L(Q),&S(PQ)&=S(P)\vee S(Q),\\
 L(P^p)&=pL(P),&S(P^p)&=
 \begin{cases}\varnothing,&p=0,\\S(P),&p>0.\end{cases}
\end{align*}
The summary of an empty cone is empty, not an empty block.

\begin{proposition}[Final-state and prefix compilation]
\label{prop:prefix}
Every summary has at most $m$ atom entries. All node summaries and expanded
lengths can be computed in polynomial time in the explicit grammar, $m$,
and the exponent bit lengths. For a supplied $0\le t\le L(P)$, the summary
of the first $t$ expanded instructions can be computed by visiting at most
one descending grammar path and joining at most $g$ stored summaries.
\end{proposition}

\begin{proof}
Each active subpartition consists of disjoint blocks, so its total atom
entries are at most $m$. The recursive summary rules are correct by
\cref{lem:semilattice}; backward references make their bottom-up evaluation
well-defined.

For a prefix of $PQ$, descend into $P$ if $t\le L(P)$. Otherwise include
$S(P)$ and descend into the first $t-L(P)$ instructions of $Q$. For a
prefix of $P^p$, descend into $P$ if $t<L(P)$. If $t\ge L(P)>0$, the prefix
already contains one full copy of $P$. Its remaining partial or complete
copies add only relations already present in $S(P)$, so its entire summary
is $S(P)$. A full-node prefix returns its stored summary immediately.
Every descent follows a strict backward reference and therefore has length
at most $g$. Zero-length nodes are handled by their empty summary and never
cause division by zero.
\end{proof}

With union by rank and path compression, normalizing and joining explicit
blocks costs $O(m\log(m+1))$ index operations using the supplied canonical
sorting conventions. If $a$ is the number of atom occurrences in the
explicit leaves, an adequate compilation bound is
\[
 O\bigl((a+gm)\log(m+1)\bigr)
 \quad\text{index operations, plus }O(g\mathsf M(H))\text{ bit operations}.
\]
Indeed $H\le g+\sum\bits(p)$ over the power nodes is a safe upper bound:
a concatenation adds at most one bit on a dependency path and a power adds
at most the bit length of its exponent. The input itself and the output
positions therefore account for all large integers.

\subsection{Direct first-threshold search}

Write $C(t)$ for the component count after the first $t$ instructions.
Since cones only coarsen partitions, $C(t)$ is nonincreasing. Prefix
bisection would locate the first $t$ with $C(t)\le K$ in
$O(\log(L(P)+1))$ exact prefix evaluations. This is already polynomial in
$H$, but its evaluation count can be much larger than the grammar size.

\begin{theorem}[Grammar-directed first threshold]
\label{thm:threshold}
Given a compiled profile table and a $g$-node fixed-port program, the first
prefix whose component count is at most an integer $K$ can be found with
at most $g+1$ quotient-state trials. The answer is $0$ if the initial count
qualifies and is \texttt{None} if the final count does not qualify.
The trial count is independent of the numerical expanded length; arithmetic
on its $H$-bit indices is still charged.
\end{theorem}

\begin{proof}
First test the initial count and, if necessary, the root's complete summary.
If neither provides a terminal answer, maintain a current quotient state,
a node, and an expanded offset such that the target is not yet reached but
will be reached somewhere inside that node.

At a concatenation $PQ$, apply $S(P)$ to a clone of the current state. If
its count qualifies, descend into $P$ without changing the original state.
Otherwise adopt the trial state, add $L(P)$ to the offset, and descend into
$Q$. At a positive power $P^p$, descend into its first copy of $P$: after
that copy the state already has the full effect of all $p$ copies by
idempotence. Thus a target attained by the power is attained in its first
copy. A zero power cannot satisfy the maintained invariant. At a cone
leaf, apply the cone and return the offset plus one.

The invariant is preserved in every case. Each step follows a strictly
backward grammar reference, so there are at most $g$ visited nodes. Only
concatenations need a new trial, besides the initial full-root test. This
gives the stated bound. Monotonicity and the left-first descent prove that
the returned prefix is the first one, not merely one qualifying prefix.
\end{proof}

\begin{example}[A huge delayed threshold]
Let $q$ parallel components meet each of eight atoms once. Let $I$ be the
empty cone and $A=\Cone(\{3\})$. The five-node program
\[
 P=(I^{2^b}A)^{2^b}
\]
has expanded length $2^{2b}+2^b$. Its first count at most one occurs at
$t=2^b+1$. The direct algorithm needs two state trials for every $b$.
The included prefix-bisection baseline uses $2b+2$ evaluations on this
family. Reading $b+1$ exponent bits and writing the event position remain
necessary; the theorem eliminates repeated state evaluation, not bit
complexity itself.
\end{example}

\subsection{Enumerating every effective event}

\begin{theorem}[Complete event-driven history]
\label{thm:history}
For a $g$-node fixed-port program, all expanded indices at which its
component partition changes can be enumerated without expanding repetitions.
If there are $h$ such events, the algorithm uses $O((h+1)g)$ quotient-state
trials and returns $h$ records of the form $(t,C(t-1),C(t))$.
In particular $h\le\max(0,2m-1)$ and the entire history is polynomial in
the encoded input parameters.
\end{theorem}

\begin{proof}
Traverse the grammar from the root with a current quotient state and an
expanded offset. At each encountered node, tentatively apply its complete
summary. If the count is unchanged, skip the entire node. This is sound
not just for counts: a coarsening of a finite partition with the same number
of blocks is the identical partition. Hence all its intermediate states
are also unchanged.

If the count decreases, a cone leaf produces an event and its trial state
is adopted. A concatenation is explored left to right, with the right
offset increased by the full expanded left length. A positive power is
explored only in its first child copy. After that copy no later copy can
change the relation, by idempotence.

Every successful nonleaf trial lies on a descending path to at least one
reported event. Charge it to the first such event in its explored subtree.
Each path has at most $g$ nodes. Each successful concatenation generates at
most two child trials, so the skipped siblings add at most a constant
factor. If there are no events, one root trial suffices. This proves
$O((h+1)g)$. Correct offsets locate events in the original expanded program,
including all skipped copies. Finally apply \cref{thm:height}.
\end{proof}

The implementation uses an explicit stack, so long grammar dependency
chains do not consume Python recursion depth. The tests include 2,500
nodes and event indices beyond ordinary decimal-serialization limits.
Hexadecimal transport keeps those limits unchanged rather than modifying
interpreter-global settings.

The history theorem is stronger than a final-state shortcut. It recovers
all genuine component changes and their exact positions in a potentially
enormous schedule. It does not recover an expanded list of redundant
operations, whose output size would itself forbid a polynomial bound.
